\section{数学の言葉を作る}

この本では、制御工学を理解するために必要な数学を本文の外へ追い出さない。読者がまだベクトル、行列、微分積分を学んでいないとしても、本文で使う範囲はここから順に作る。数学を先に大量に準備してから制御へ入るのではなく、制御で必要になる言葉を、使う直前に意味と式の両方で確認する。

\begin{definition}[量と単位]
数値と単位を合わせたものを量という。たとえば $3\,\mathrm{s}$ は時間、$5\,\mathrm{m}$ は長さ、$12\,\mathrm{V}$ は電圧を表す量である。
\end{definition}

数だけを見ても、物理的な意味は決まらない。$10$ が $10\,\mathrm{m}$ なのか、$10\,\mathrm{s}$ なのか、$10\,\mathrm{rad/s}$ なのかでまったく意味が違う。制御工学では式を作るとき、左辺と右辺の単位が一致していなければならない。たとえば位置 $q$ の単位が m、速度 $v$ の単位が m/s なら、
\begin{equation}
  q+v
\end{equation}
は m と m/s を足しているので、そのままでは意味を持たない。一方、短い時間 $\Delta t$ の間に速度 $v$ で進む距離は $v\Delta t$ であり、単位は
\begin{equation}
  \frac{\mathrm{m}}{\mathrm{s}}\cdot\mathrm{s}=\mathrm{m}
\end{equation}
だから、$q+v\Delta t$ は位置として意味を持つ。

\begin{lecturepoint}
式を見て迷ったら、まず単位を見る。単位は、式が物理として成立しているかを最初に検査する道具である。
\end{lecturepoint}

\subsection{変数と関数}

\begin{definition}[変数]
値が状況によって変わる記号を変数という。時刻を $t$、位置を $q$、速度を $v$、入力を $u$ と書くとき、それらは変数である。
\end{definition}

\begin{definition}[関数]
ある値を入れると別の値が決まる対応を関数という。時刻 $t$ を入れると位置 $q(t)$ が決まるとき、$q$ は時刻の関数である。
\end{definition}

関数 $q(t)$ は、表でも、グラフでも、式でも表せる。たとえば一定速度 $v_0$ で動く物体の位置は
\begin{equation}
  q(t)=q(0)+v_0t
\end{equation}
と書ける。ここで $q(0)$ は時刻 0 の位置である。$t=2$ を代入すれば
\begin{equation}
  q(2)=q(0)+2v_0
\end{equation}
となる。このように、関数とは「式の形を持つもの」だけではなく、「入力を与えると出力が決まる対応」である。

\subsection{差分と平均変化率}

微分を知らなくても、変化を比べることはできる。時刻 $t_1$ の位置が $q(t_1)$、時刻 $t_2$ の位置が $q(t_2)$ なら、位置の変化量は
\begin{equation}
  \Delta q=q(t_2)-q(t_1)
\end{equation}
である。時刻の変化量は
\begin{equation}
  \Delta t=t_2-t_1
\end{equation}
である。

\begin{definition}[平均変化率]
量 $q$ が時間 $\Delta t$ の間に $\Delta q$ だけ変化したとき、
\begin{equation}
  \frac{\Delta q}{\Delta t}
\end{equation}
を平均変化率という。
\end{definition}

たとえば $t_1=1\,\mathrm{s}$ で $q(t_1)=2\,\mathrm{m}$、$t_2=4\,\mathrm{s}$ で $q(t_2)=11\,\mathrm{m}$ なら、
\begin{align}
  \Delta q&=11-2=9\,\mathrm{m},\\
  \Delta t&=4-1=3\,\mathrm{s},\\
  \frac{\Delta q}{\Delta t}&=\frac{9\,\mathrm{m}}{3\,\mathrm{s}}=3\,\mathrm{m/s}.
\end{align}
これは「この区間では平均して毎秒 3 m 進んだ」という意味である。

\begin{why}
微分は、平均変化率の時間幅 $\Delta t$ を限りなく小さくしたものとして現れる。したがって微分を学ぶ前に、差分と平均変化率を理解しておけば、微分は突然の記号ではなく、すでに知っている考えの極限になる。
\end{why}

\subsection{複数の数を並べる}

制御対象には、一つの数だけでは表せないものが多い。平面上の位置には横方向と縦方向があり、モータには電流、角速度、角度がある。このように複数の数を順番に並べたものを、後でベクトルと呼ぶ。

\begin{definition}[ベクトル]
複数の数を順番に並べたものをベクトルという。たとえば
\begin{equation}
  x=\begin{bmatrix}2\\3\end{bmatrix}
\end{equation}
は、1番目の成分が 2、2番目の成分が 3 のベクトルである。
\end{definition}

ベクトルは「複数の情報を一つの記号で持つ」ための入れ物である。平面上の位置なら
\begin{equation}
  p=\begin{bmatrix}p_x\\p_y\end{bmatrix}
\end{equation}
と書ける。速度も同じように
\begin{equation}
  v=\begin{bmatrix}v_x\\v_y\end{bmatrix}
\end{equation}
と書ける。位置と速度をまとめて状態にするなら
\begin{equation}
  x=\begin{bmatrix}p_x\\p_y\\v_x\\v_y\end{bmatrix}
\end{equation}
と並べればよい。

\begin{definition}[行列]
数を長方形に並べたものを行列という。たとえば
\begin{equation}
  A=\begin{bmatrix}1&2\\3&4\end{bmatrix}
\end{equation}
は 2 行 2 列の行列である。
\end{definition}

行列は、ベクトルを別のベクトルへ変える規則として使える。たとえば
\begin{equation}
  A=\begin{bmatrix}1&2\\3&4\end{bmatrix},
  \qquad
  x=\begin{bmatrix}5\\6\end{bmatrix}
\end{equation}
のとき、行列とベクトルの積は
\begin{align}
  Ax
  &=\begin{bmatrix}1&2\\3&4\end{bmatrix}
    \begin{bmatrix}5\\6\end{bmatrix}\\
  &=\begin{bmatrix}1\cdot5+2\cdot6\\3\cdot5+4\cdot6\end{bmatrix}\\
  &=\begin{bmatrix}17\\39\end{bmatrix}.
\end{align}
この計算を省略せずに見れば、行列は「謎の記号」ではなく、各行がベクトルの成分を重み付きで足し合わせる表だと分かる。

\subsection{添字と総和記号}

\begin{definition}[添字]
同じ種類の量を区別するために、文字の右下につける番号を添字という。$x_1,x_2,x_3$ は、ベクトル $x$ の 1 番目、2 番目、3 番目の成分を表す。
\end{definition}

ベクトル
\begin{equation}
  x=\begin{bmatrix}x_1\\x_2\\x_3\end{bmatrix}
\end{equation}
の成分をすべて足すと
\begin{equation}
  x_1+x_2+x_3
\end{equation}
である。成分がたくさんあると毎回書くのが長いので、総和記号を使う。

\begin{definition}[総和記号]
記号
\begin{equation}
  \sum_{i=1}^{n}x_i
\end{equation}
は、$i$ を 1 から $n$ まで動かして $x_i$ を足すことを表す。
\end{definition}

たとえば $n=3$ なら
\begin{equation}
  \sum_{i=1}^{3}x_i=x_1+x_2+x_3
\end{equation}
である。内積やノルム、評価関数、確率の期待値は、この総和の考え方を使って短く書かれる。ここまでの記号が分かれば、次に出てくる信号、微分、状態方程式を「記号の壁」としてではなく、量の変化を整理する言葉として読める。